\section{The normal distribution}
\label{sec:normal-dist}

We say that $X$ has the normal (a.k.a. gaussian)
distribution with mean $\mu$ and variance $\sigma^2$,
and denote $X\sim N(\mu,\sigma^2)$, if $X$ has
density function

\[
f_X(x) = \frac{1}{\sqrt{2\pi} \sigma} e^{-(x-\mu)^2/2\sigma^2} \qquad (x\in\R).
\]

\noindent In particular, the standard normal distribution is
the distribution $N(0,1)$, whose density function is
given by

\[
f_X(x) = \frac{1}{\sqrt{2\pi}} e^{-x^2/2} \qquad (x\in\R).
\]

\noindent $\textbf{Properties:}$

\begin{enumerate}[label=\arabic*.]
\item If $X\sim N(\mu,\sigma^2)$ then $\mathbf{E}(X)=\mu$,
$\mathbf{V}(X) =\sigma^2$.

\item $N(\mu_1,\sigma_1^2) \boxplus N(\mu_2,\sigma_2^2) =
N(\mu_1+\mu_2,\sigma_1^2+\sigma_2^2)$.

\item If $X,Y\sim N(0,1)$ are independent and standard
normal then $\frac{1}{\sqrt{2}}(X+Y) \sim N(0,1)$.

\item More generally, if $X_1,\ldots,X_n \sim N(0,1)$ are
independent standard normal r.v.s then

\[
\frac{1}{\sqrt{n}}(X_1+\ldots+X_n) \sim N(0,1),
\]

\noindent and also

\[
\sum_{j=1}^n \alpha_j X_j \sim N(0,1)
\]

\noindent if $\alpha_1,\ldots,\alpha_n$ are real numbers such
that $\sum_j \alpha_j^2 = 1$. (Geometrically,
$\mathbf{\alpha}\cdot \mathbf{X} = \sum_{j=1}^n
\alpha_j X_j$
can be interpreted as the projection of the random
vector $(X_1,\ldots,X_n)$ in $\R^n$ in the direction
of the unit vector
$\mathbf{\alpha}=(\alpha_1,\ldots,\alpha_n)$.)

\item If $X\sim N(0,1)$ then

\[
\mathbf{E}(X^k) = \begin{cases}
0 & \textrm{if } k \textrm{ is odd}, \\
1\cdot 3 \cdot 5 \cdot \ldots \cdot (k-1) & \textrm{if } k \textrm{ is even}.
\end{cases}
\]

\item The normal distribution is the single most important
distribution in probability! The theoretical reason
for this is the $\textbf{Central Limit Theorem}$, a
result we will discuss in detail in chapters 11--14.

\item The polar decomposition of a bivariate standard
normal vector: given a pair $(X,Y)$ of random
variables which in the polar representation are
written as $X=R\cos\Theta$, $Y=R\sin \Theta$, where
$R>0$ and $0\le \Theta<2\pi$, we have

\begin{align*}
X,Y\sim N(0,1), & \ X,Y\textrm{ are independent } \\
& \iff R^2 \sim \operatorname{Exp}(1/2), \ \Theta\sim U[0,2\pi]\ \textrm{and } R,\Theta \textrm{ are independent.}
\end{align*}
\end{enumerate}
